% !TeX root = ../papers/nc-p3.tex
\section{単元座標での極の全深度判定}
\label{sec:rank-fifteen-unit-pole-depth}
\begin{theorem}\label{thm:rank-fifteen-unit-pole-depth}
標数零の通常な階数15窓で、剰余標数3の付値に対して $\nu(x)=0$, $\nu(y)>0$ とする。
次のいずれかの場合には、付値の深さによらず標準Koszul左右AS正則であり、
各始点で四中心四次式上自由階数256である。
\begin{enumerate}
\item $\bar x\ne1$.
\item $\bar x=1$, $2\nu(x-1)\ge\nu(y)$ で $J=(x-1)^2/(4y)$ の剰余が1と異なる。
\end{enumerate}
\end{theorem}
\begin{proof}
\path{computations/verify_rank_fifteen_unit_pole_depth.py} は各表示に十七桁を取り、
$\mathbb Q(s)$ または $\mathbb Q(z)$ 上で、実際に使う全てのパラメータ除算の前に
定数項が恒等的に零であることを確認する。係数環はそれぞれ
$\mathbb Z[s]_{(3,s^4+2s^3+s+1)}$、$\mathbb Z[z]_{(3,z^2+z+2)}$ である。
主元の定数項の分子・分母の3での還元を記録する。

第一の場合、未切断経路と中心方程式、五次の中心商に九桁を使い、
長さ6から13の最大損失は $(2,1,1,1,1,0,1,0)$ である。
従って形式級数環上の語の格子の両積表を最後まで持ち上げられ、十三次で商は零となる。
使用した定数係数の平方自由な積は、同じ深さの $\mathbb F_3(s)$ 上の
計算で得た次数121の多項式に一致する。第\ref{sec:rank-fifteen-pole-fibers}節の
八重被覆で残る二つの三次式のファイバーは追加表示で覆う。

この比較では、基準と追加表示の全308行列のSmith指数の重複度を照合する。
基準の標数零の計算では単元主元が尽きたときに分数体でも定数項が零なので、
基準の剰余体での指数列が分数体上の指数列に一致する。
補題\ref{lem:rank-fifteen-smith-transfer}を順に適用すると、
追加表示でも同じ本数の主元が混標数の格子上で得られる。
中心四次式を指定する最後の始点の語座標は同じであり、
左右からの積と中間語の基底変更は照合行列の可逆な行・列変換になる。
語の格子の自由性は一般点の次元下界による。従って基準と追加表示のいずれでも、
正付値のパラメータへの特殊化後に同じ上界を与える。

第二の場合、基準表示で十二次までの292行列を標数零で消去する。
五次までに十桁、続く長さ6から12では $(1,0,1,1,2,0,1)$ 桁を使い、
一桁を残す。全ての指数列と中心座標が剰余のものに一致する。
ここでの主元の平方自由な積は、第\ref{sec:rank-fifteen-ramified-pole}節で
十三次までを含めて得た次数897の多項式と等しい。
同節の五つの追加表示でこれらの指数列を移すと、十二次までの同じ自由格子を得る。
十三次は各表示でパラメータ除算を使わず単元主元だけで全列を消す。
従って十二次で特殊化した積表の剰余から十三次の消滅を持ち上げられる。
両方の場合に $(1+t+t^2+t^3)^4$ の全係数と商の上界が一致し、
第\ref{sec:rank-fifteen-overlap-annulus}節の左右の解消と自由加群の議論が適用できる。
\end{proof}
\begin{corollary}
$x,y\in\mathbb Q$, $y\ne0$ が通常な階数15窓を与えるとする。
$x$ が3進単元で $x\equiv-1\pmod3$ なら、全ての $y$ で上の正則性と自由性が成り立つ。
$x\equiv1\pmod3$ でも $0<\nu_3(y)<2\nu_3(x-1)$ の全点が含まれる。
特に $\nu_3(x)=0$, $\nu_3(y)=1$ の全有理点が含まれる。
\end{corollary}
\begin{proof}
負付値の $y$ には極の定理、単元 $y$ には素体単元の四点の定理が適用できる。
正付値では第一の場合か、厳密な不等式により $\bar J=0$ となる第二の場合を使う。
\end{proof}
